\documentclass[11pt]{article}
% Shared preamble for all daily exercises
% Update this file to change settings (padding, parindent, etc.) for all exercises

\usepackage[margin=0.75in]{geometry}
\usepackage{amsmath}
\usepackage{amssymb}

% Custom settings
\setlength{\parindent}{0pt}
\setlength{\parskip}{0.2cm}

\title{Daily Exercise}
\author{Dhruman Gupta}
\date{24th April}

\begin{document}

% Common header for daily exercises
% This file can be included in other LaTeX documents

\maketitle

\noindent
\textbf{Course:} CS-3330, Theory of Computation

\vspace{0.2cm}

\noindent
\textbf{Question:}\noindent 
Let $f$ be a time-constructible function, that is, $f(n)$ can be computed in $O(f(n))$ time and $f(n) \geq n$.

Show that if $NTIME(n) = DTIME(n)$, then $NTIME(f(n)) = DTIME(f(n))$. Is the converse true?

\textbf{Answer:}

Assume that

\[
NTIME(n) = DTIME(n).
\]

We want to show that

\[
NTIME(f(n)) = DTIME(f(n)).
\]

Since $DTIME(f(n)) \subseteq NTIME(f(n))$, it is enough to prove that

\[
NTIME(f(n)) \subseteq DTIME(f(n)).
\]

Let $L \in NTIME(f(n))$. Then some nondeterministic TM $M$ decides $L$ in time $O(f(n))$.

Now define the padded language

\[
L' = \{x\#^{f(|x|)-|x|} \mid x \in L\}.
\]

Since $f(n) \geq n$, this is well-defined, and every string in $L'$ has length exactly $f(|x|)$.

We claim that $L' \in NTIME(n)$. Given an input $y$, a nondeterministic machine first checks whether $y$ has the form $x\#^t$. It then computes $f(|x|)$ in $O(f(|x|))$ time, checks that

\[
|y| = f(|x|),
\]

and then simulates $M$ on $x$.

All of this takes $O(f(|x|))$ time. But since $|y| = f(|x|)$, this is just $O(|y|)$. Hence,

\[
L' \in NTIME(n).
\]

By assumption $NTIME(n) = DTIME(n)$, so

\[
L' \in DTIME(n).
\]

Now use this to decide $L$. Given input $x$, compute $f(|x|)$, construct the padded string

\[
x\#^{f(|x|)-|x|},
\]

and run the deterministic linear-time machine for $L'$ on this string.

The constructed string has length $f(|x|)$, so the simulation takes $O(f(|x|))$ time. Computing $f(|x|)$ and writing down the padding also take $O(f(|x|))$ time. Therefore $L$ is decidable deterministically in time $O(f(n))$.

So,

\[
L \in DTIME(f(n)).
\]

Since $L$ was arbitrary, we get

\[
NTIME(f(n)) \subseteq DTIME(f(n)).
\]

Hence,

\[
NTIME(f(n)) = DTIME(f(n)).
\]

Now for the converse. The converse would say that if

\[
NTIME(f(n)) = DTIME(f(n)),
\]

then

\[
NTIME(n) = DTIME(n).
\]

This is not true in general. Equality at a larger time bound does not imply equality at a smaller time bound. For example, knowing that nondeterminism adds no power at time $n^2$ does not immediately imply that it adds no power at time $n$.

So the converse is false in general.

Hence proved.

\end{document}
